% ── §9 Methods: Validation Framework ──
\section{Validation Framework}
\label{sec:methods-validation}

\subsection{Statistical Methods}
All enrichment tests use exact one-sided binomial tests (scipy
\texttt{binomtest}) against matched-separation random controls constructed by
drawing 2{,}000 random residue pairs per model from the same protein's
candidate universe, stratified to preserve the separation-bin distribution of
the tested pair set.  Spearman rank correlations are computed using scipy
\texttt{spearmanr}.  Multiple-comparison awareness is maintained through
median (rather than mean) reporting for skewed precision distributions and
through explicit Bonferroni-corrected significance thresholds where noted.

\subsection{Determinism}
Every experiment uses fixed random seeds (protein folds: seed~42; bootstrap
resampling: seeded \texttt{default\_rng}; background construction: per-model
seeds).  Re-running any stage produces byte-identical output files, verified
by MD5 checksums across at least two independent invocations for each of the
seven major result categories.

\subsection{Code Quality}
All production scripts pass pyflakes with zero warnings.  The adversarial
unit test suite comprises 26 test functions covering encoding helpers,
alignment loading (including ragged and empty inputs), PDB parsing (including
alternate-location indicators and inconsistent residue naming), native-contact
geometry, identity-mapping verification, mutual-information equivalence with
the shared module, truth-table labeling, Quine--McCluskey soundness and
completeness on real data, implicant-decoder predicate equivalence over
randomised inputs, metric correctness (precision@k, MCC, AUC vs.\ brute-force
concordance), cross-validation fold determinism, and adversarial empty-input
handling.
